\begingroup
\setlength{\LTcapwidth}{\textwidth}
\begin{longtable}{@{}p{2.2cm}p{3.9cm}p{6.8cm}@{}}
\caption[Keempat unsur kerangka komunikasi klinis SBAR.]{Keempat unsur kerangka komunikasi klinis SBAR (Sumber: Müller dkk. 2018).}
\label{tab:tab2.2_sbar} \\
\toprule
\textbf{Unsur} & \textbf{Kepanjangan} & \textbf{Isi yang disampaikan} \\
\midrule
\endfirsthead

\multicolumn{3}{@{}l}{\small\emph{Tabel \thetable{} (lanjutan)}} \\
\toprule
\textbf{Unsur} & \textbf{Kepanjangan} & \textbf{Isi yang disampaikan} \\
\midrule
\endhead

\bottomrule
\endlastfoot

S & \textit{Situation} &
Keadaan yang sedang dihadapi serta alasan komunikasi dilakukan, disampaikan secara ringkas pada bagian awal. \\[4pt]

B & \textit{Background} &
Riwayat dan konteks klinis yang diperlukan untuk memahami keadaan, termasuk data relevan dan tindakan yang telah dilakukan. \\[4pt]

A & \textit{Assessment} &
Penilaian atas keadaan yang dihadapi, termasuk interpretasi atau dugaan penyebab serta tingkat kegentingannya. \\[4pt]

R & \textit{Recommendation} &
Tindakan yang diusulkan atau hal yang perlu dilakukan oleh penerima informasi, dinyatakan secara spesifik agar dapat ditindaklanjuti. \\
\end{longtable}
\endgroup
